\section{Personalization of Proportional Representation}

\subsection{No Independent Candidates}

The basic idea behind proportional representation is that all socially
relevant intellectual currents and social groups should be represented in parliament. The
representation takes place through parties, which were originally conceived as groups of like-minded individuals
regarding certain ideas. However, it may also happen that two parties
differ hardly in content, but are simply two teams that promise to
do the same thing better. They are comparable to two telecommunications companies that
compete for customer favor by trying to be better than the competition in the same matter, telecommunications.
There is nothing inherently wrong with this, as long as it is still
ensured that all relevant social currents are still represented in parliament by other parties if necessary.

In principle, a single person could also represent an entire intellectual current.
However, I would still not allow individuals to run for state and
federal elections for two reasons:
\begin{enumerate}
  \item \textbf{Transparency.} The parliaments in question here are parliaments of large
    units, such as entire countries. For such parliaments, it is actually inconceivable
    that an individual person could acquire the necessary expertise and recognition without an organization
    backing them up. If someone wants to run as an independent candidate, then their
    organization should also be made transparent by having the candidate found their own party.

  \item \textbf{Procedural reasons.} Both the iterative electoral procedure presented above and
    the constituency representation procedure, which will be introduced later in this document,
    work at the party level. Depending on the vote result, a party receives more or fewer
    seats in parliament and represents more or fewer constituencies. How should a single candidate
    receive more or fewer seats in parliament? Independent candidates would disrupt the entire procedure.
\end{enumerate}

\runinhead{Conclusion:}

A selection of persons by the voter must be made within parties. The
voter must be offered an opportunity to make a selection among several candidates of the same party.

\runinhead{Note:}

The prohibition of independent candidates does not imply that there could not be independent representatives.
If during a legislative period there is a falling out between a representative
and the party that sent them, this representative may continue to exercise their mandate as a now
independent representative. This permission affects neither transparency nor the
procedure, but is absolutely necessary to ensure that a representative's loyalty
can be first to their conscience and the voter, and only secondarily to party interests.
If this order were reversed, parliamentary democracy would degenerate into
dictatorship.

\subsection{Election of Lead Candidates by the People}

Especially with the large, established parties, there is the phenomenon that within the
parties, people advance in their careers who are not liked by the voters. This is no accident. To
advance in a party, one must please or prevail with a different group of people
than in a parliamentary election. As a result, voters often face
the decision of whether to vote for a party despite or because of its lead candidate. Often
one agrees with a party on content, but cannot identify with their
lead candidate. The CDU lost the 2021 federal election for this reason. Conversely,
it can of course also be the case. One likes the candidate, but cannot support the positions of their party.

To mitigate the problem, I propose that all parties running for election must
run with three lead candidates. Each voter receives a \enquote{party vote} with
which they determine the majority ratios in parliament, and a
\enquote{lead candidate vote} with which they can choose among the lead candidates of their party.
The ballot shows a table with 4 columns: The first column lists the
parties. Next to them are the three lead candidates of each party. The voter makes their
mark for a party and can also make a mark for one of the candidates of the party they have chosen.
A mark for a candidate of another party is not permitted. The party wins
(possibly after several iterations) based on the party votes, and the lead candidate
of this party who was able to garner the most votes.

\subsection{Procedure for Direct Election of Representatives Without Changing the Party Composition}

In the current German electoral law, an influence on the personal composition of the
Bundestag by the voter is provided. With their \enquote{first vote}, one could previously promote a
candidate to the Bundestag via direct mandate independent of the second vote result. When
first and second vote results differed, compensatory mandates arose to maintain
proportionality, which led to an enlargement of the Bundestag. Under the new electoral law,
first vote wins only count to the extent that they are covered by the second vote result. Mandates
due by second vote are first filled with constituency winners (if available) in the order
of their relative vote shares, before continuing with state lists. This allows
voters to virtually place candidates at the front of the state list and thus theoretically influence the personal
composition of the Bundestag factions. In practice, however, this only results in
changes if the parties don't already place their constituency candidates at the front of the lists anyway.

The old procedure was mainly criticized because the Bundestag was inflated by overhang and
compensatory mandates when voters made use of their right to split voting.
This problem can no longer occur with the new electoral law. The
Bundestag always has exactly 630 seats.

The new electoral law was criticized because there can be constituencies that are not represented in the Bundestag,
as they were won by party candidates whose parties did not have enough
second vote shares. This would not have happened under the old electoral law, as a
constituency win meant a secure place in the Bundestag.

To see whether anything can be improved here, one must first take a step
back and consider what one actually wants to achieve and can achieve with this \enquote{representative election law}.

\begin{enumerate}
  \item One wants representatives to sit in parliament who are personally known
    by voters and accountable to them. This increases both the legitimacy and the
    control of parliament by the citizens. If one wants to be
    reelected in their constituency, one cannot promise the moon and then
    do the opposite. At least in rural areas, voters know their representatives,
    and will not accidentally elect criminals to parliament. If they elect criminals to
    parliament, they do so intentionally because they share their convictions.

  \item One wants all areas on the sovereign territory of parliament to be represented
    in parliament as well.
\end{enumerate}

The first point can be interpreted with varying degrees of strictness. One can a) demand that
enough representatives personally known by citizens sit in parliament so that control of
parliament is ensured. Or one demands b) that all representatives of parliament are
personally known by citizens. I consider demand~b) to be too strong, as parliament should also
possess substantive competence. While citizens can and must assess trustworthiness very well,
this is generally not the case with substantive competence. The parties must
do this themselves, as they know their experts precisely. An expert must be good in their field,
but not necessarily in citizen communication. Therefore, I would suggest being content with
demand~1a) and only filling half of parliament with candidates personally known by citizens.
(This is also roughly the case in the current electoral law.)

The second point can generally only be fulfilled with nationwide proportionality and constant parliament size if one understands
\enquote{representation of areas} not as
\enquote{representation of local power}.
The nationwide\footnote{\enquote{Nationwide} or \enquote{country} refers hereinafter to
the sovereign territory of the respective parliament. For a federal election, this would be the
federal territory.} power ratio of the parties generally does not correspond to the local distribution
of the parties. If, for example, party~A unites 60\,\% of the nationwide voter votes,
party~B 30\,\% and party~C 10\,\%, this normally does not mean that in 60\,\% of the constituencies
party~A is the strongest force, in 30\,\% of the constituencies party~B and in 10\,\% of the constituencies
party~C. If all parties run nationwide and really represent the interests of the entire country,
one should expect party~A to become the strongest force in every constituency. This
is precisely the strength of proportionality, as even with such strong nationwide superiority of one
party, the others are still represented in parliament. Experience shows that even nationwide
competing parties often have regional strongholds where they traditionally
perform strongly in elections and their political opponents correspondingly weaker. Therefore, it could be that
even in the above example, party~B or~C still win one or another constituency.
It would be problematic for the cohesion of a country, however, if indeed 60\,\% of all
constituencies vote for party~A, 30\,\% of all constituencies for party~B and 10\,\% of all constituencies for party~C,
and this even in contiguous areas. This would mean that obviously no
party has the interests of the entire country in view, but only their respective core area.
Moreover, area A would virtually rule over areas B and C. This provides fuel for
separatist efforts or even civil wars.
In this sense, the result of the 2025 federal election is very problematic, where the entire East voted blue,
and the West black with some red enclaves.\footnote{\url{https://www.tagesschau.de/inland/bundestagswahl/wahlkreisergebnisse-karte-bundestagswahl-2025-100.html}} This is a
trend that must not continue!
However, if we assume a \enquote{healthy} regional vote distribution, i.e., the
power ratios between the parties are roughly the same throughout the country apart from a few
regional fluctuations (\enquote{strongholds} or coincidences), we must also fill parliamentary seats of
parties with candidates who have won in no constituency and therefore
take away their constituency from \enquote{legitimate constituency winners}. This error is
inherent to the principle. It cannot be avoided, only minimized. Somewhere we must
shift the error: Either we make the size of parliament variable, or we live with the fact that
some constituencies have no representation in parliament, or we take away the constituency from some
constituency winners and give them to a loser, so that the seat ratio in
parliament corresponds to the nationwide voter will. Parliament enlargements have long been
unpopular, missing representation of some constituencies was criticized in the 2025 federal election.
Only one option remains: taking away the constituency from some constituency winners again.
A departure from nationwide proportionality would also be a possibility. If the
party composition of parliament were determined by party affiliation of constituency winners,
every constituency would be represented, but the majority ratios in parliament would
no longer represent the nationwide power ratio of the parties. An iterative
procedure, such as we have also proposed, would then no longer make sense, as it
would change the majority ratios again.

The way out consists in only demanding that every area has representation in the
nationwide power and not that the local power ratios must be
reflected in parliament.

\enquote{Representation in nationwide power} can be interpreted with varying degrees of strictness:
\begin{enumerate}[label=\alph*)]
  \item Each constituency has a representative in the governing party or
    governing coalition.
  \item Each constituency sends a representative to parliament. However, this does not have to be a
    representative in the governing coalition. Some constituencies must be content with
    parliamentary representation through opposition politicians.
\end{enumerate}

Demand~a) is stronger and more just than demand~b), as access to the
governing party/coalition naturally brings potentially more opportunities for influence, and
constituencies without government representation, which would exist under~b), would be
disadvantaged compared to those with government representation. However, demand~a) also requires approximately double the
size of parliaments as demand~b).
We had demanded above for good reason that only half of parliament be directly
elected by citizens. The other half should be filled by the parties so that they can
contribute their experts, who do not necessarily have to be good at citizen communication. If
demand~a) applies, then this principle must also apply to the governing party/coalition,
which generally only has slightly over half of the parliamentary seats. If it is guaranteed that all
constituencies always have government representation, then the number of constituencies must not
exceed one quarter of the parliamentary seats. However, the constituencies should have a size
such that it is still possible to know representatives personally. According to the author's
experience, this is actually still the case with 200,000 -- 300,000 inhabitants per constituency.
From the maximum number of eligible voters per constituency and the number of
eligible voters in total, however, the minimum number of constituencies results, from which
the minimum parliament size results as a) $4\times$ minimum number of constituencies and in
variant~b) as $2\times$ minimum number of constituencies.

In Germany, we currently have 299 constituencies, which average about 240,000 inhabitants.\footnote{\url{https://www.bundestag.de/parlament/bundestagswahl/wahlkreise-1050626}} With
630 parliamentary seats, variant~a) would therefore not be feasible. Now one could increase both the
number of inhabitants per constituency and the parliament size to implement variant~a),
or one accepts the unequal treatment of constituencies arising from variant~b)
in favor of a smaller parliament.

The government of a territory must absolutely act in the interest of the entire
country in all three branches of government and, if necessary, make trade-offs between conflicting regional interests.
Even if constituency representatives do bring a kind of regional representation back into play through the back door,
this should play a subordinate role, as
in a national parliament\footnote{not \enquote{federal state parliament}}, such as the German
Bundestag, politics is explicitly made and laws are passed that apply to the entire country.
Furthermore, there was also no claim under the previous system of direct mandates
to area representation by a governing party/coalition representative. Rather, one could
elect one's representative to parliament. Whether they then belonged to the governing majority or not
depended on the power ratios in the Bundestag. No one has ever been bothered by the resulting
area inequality. Therefore, I think one could live with variant~b).
In the end, however, it is a societal decision how much equal treatment
of areas is worth to us. Such questions could possibly also be put to the people for a vote.

Regardless of which variant one chooses, the problem remains of determining the local
constituency representative without changing the party composition in parliament. The
proposal is that each party must run with at least three candidates in each constituency.
Each voter can designate a candidate as their preferred representative for each party in their constituency.
They can do this for every running party, not just the one they
voted for with their \enquote{party vote}. Each voter thus receives 3 types of
votes:
\begin{enumerate}
  \item A \enquote{party vote} with which they elect a party to parliament.
  \item A \enquote{lead candidate vote} which they can cast for a lead candidate of the party they chose with~1.
  \item A \enquote{representative vote} per running party. For each running party, they can
    choose their favorite or least hated candidate in their constituency.
\end{enumerate}
After the party composition of parliament is determined, according to a mathematical
procedure that considers the relative majorities of the parties per constituency and is
supplemented by a random procedure in case of non-uniqueness, each constituency is assigned to exactly one party.
The candidate of this party who unites the most representative votes is
the representative of this constituency. It is thus not possible to change the
numerical ratio of the parties in parliament again by electing a person, after it has already been determined by the
party votes. Nevertheless, it is possible to influence the personal composition
of parliament by making use of one's representative vote even for other parties.

\subsection{Procedure for Assigning a Constituency to a Party (Constituency Representation Procedure)}

An electoral area $W$ (e.g., a country) is divided into $M$ constituencies $w_1 \ldots w_M$. Each
constituency corresponds to a parliamentary seat to be filled by direct election by the citizens according to the above procedure.
Thus, $M$ parliamentary seats are to be distributed. Since we explicitly also
allow coalitions for government formation, just no longer want to force them, it does not matter for the
procedure whether $M$ corresponds to half of parliament or only half of the governing coalition.
The procedure must be able to divide $M$ constituencies among $N$ parties. In the case that
the above variant~a) is applied and a party achieves the absolute majority or is
awarded it, $N$ is simply 1.

From the result of the party votes, a direct mandate distribution $\mu$ results, which assigns
to the parties $p_1 \ldots p_N$ the number of mandates (parliamentary seats) $m_1 \ldots m_N$.
Which procedure is used to determine the direct mandate distribution $\mu$ from the party vote ratio for a given total number $M$
plays no role for the following explanations of
how a rational vote ratio can best be approximated by an integer seat distribution.
Over time, various procedures have been developed for this problem, whose advantages and disadvantages have long been the subject of research.
As an entry point into the subject, for example, \enquote{Wahlrecht und Parteiensystem} by Dieter Nohlen may be consulted.

In the following, $m_1 \ldots m_N$ is assumed to be given.

The parties $p_1 \ldots p_N$ receive in the constituencies $w_1 \ldots w_M$ respectively the votes

\[
  \underline{S} = \begin{pmatrix}
    s_{11} & \cdots & s_{1N} \\
    \vdots & s_{ij} & \vdots \\
    s_{M1} & \cdots & s_{MN}
  \end{pmatrix}
\]

with $s_{ij}$ -- number of party votes in the $i$-th constituency for the $j$-th party.
However, what is actually interesting are only the relative vote shares per constituency

\[
  \underline{R} = \begin{pmatrix}
    r_{11} & \cdots & r_{1N} \\
    \vdots & r_{ij} & \vdots \\
    r_{M1} & \cdots & r_{MN}
  \end{pmatrix}
\]

with $r_{ij} = s_{ij}/S_i$, where $S_i$ stands for the total number of party votes cast in the
$i$-th constituency. The sum of the relative vote shares in a constituency is of course
one. Expressed in formulas:
\[
  S_i = \sum_{j=1}^{N} s_{ij}
  \qquad \text{and} \qquad
  \sum_{j=1}^{N} r_{ij} = \sum_{j=1}^{N} \frac{s_{ij}}{S_i}
    = \frac{1}{S_i}\sum_{j=1}^{N} s_{ij} = 1
\]

We seek a procedure $V(R)$ that uniquely assigns each constituency $w_i$ to a party $p_j$.
$N \leq M$ should hold, i.e., it can happen that a party is assigned to several constituencies,
but it can never happen that a party receives no constituency at all
because too few constituencies are available.\footnote{This condition will in practice probably
always be satisfied in federal elections, but is strictly speaking not necessary. If $N > M$
holds, one can always order the parties in descending order by their party vote results and
divide the constituencies among the first $N' = M$ parties. The remaining parties are then
represented in parliament, but no longer represent any constituency.}

For the case $N = 1$, the procedure is trivial. All constituencies are assigned to the only available
party.

For $N = 2$, one can sort the constituencies from the perspective of each party in descending order by their relative
vote shares. Thus, for example, for $M = 5$, the following order $\Omega_1$ could result for party $p_1$:
\[
  \Omega_1 = [w_3, w_5, w_2, w_4, w_1]
\]
This would mean that $r_{31} > r_{51} > r_{21} > r_{41} > r_{11}$ holds.

One could now determine $\Omega_2$ as a list ordered by relative election result for $p_2$ in the same way.
For $N = 2$, however, $\Omega_2$ is already implicitly given by $\Omega_1$ as $\Omega_1$ in
reverse order, since $r_{i1} + r_{i2} = 1$ holds. Thus $r_{i2} = 1 - r_{i1}$. Where
the first party got relatively the most votes, the second party got the fewest,
and vice versa.

Thus, for the above example, $\Omega_2$ would result as $\Omega_2 = [w_1, w_4, w_2, w_5, w_3]$.

These $\Omega$ lists state how important this party is to the citizens in a constituency, in
relation to the importance that citizens in other constituencies attach to it. From this, one can of course also
conversely derive an importance of the constituencies for the party: If a party
can only represent some of the constituencies (which is always the case for $N > 1$), then it should
of course represent the part where it is more wanted than in the other part. For
$N = 2$, this can always be fulfilled, as the importance is exactly reversed. For
$N = 2$, one can simply assign each party the first $m_j$ constituencies from the $\Omega_j$ list.

Then all constituencies will be assigned, and each party represents the most important
part of the constituencies for it.

What do we do in the case $N > 2$? For $N > 2$, the same constituency can indeed be the most important for several
parties, just not for all at the same time. This can be easily visualized. The
image below shows an exemplary distribution of relative vote shares for three
parties, where the 2nd and 3rd parties both simultaneously have their maximum relative
vote shares where the first party has its minimum. (On the $y$-axis, the constituencies are
plotted. The distances between the lines or the edges are the relative vote shares
of the parties. The sum of the relative vote shares is 1. For simplicity, we assume
there are so many constituencies that we no longer see the gaps between the points and we can
thus draw the lines continuously.)

\bigskip
\begin{center}
\includegraphics[width=0.8\textwidth]{diagrams/maximum_for_2_and_3_same_constituency.pdf}
\end{center}
\bigskip

It could also have been different. In the following example, all 3 parties have their maximum
relative vote shares in different constituencies.

\bigskip
\begin{center}
\includegraphics[width=0.8\textwidth]{diagrams/maxima_in_different_constituencies.pdf}
\end{center}
\bigskip

Since we already have a working procedure for $N = 2$, one could try to reduce the
procedure for $N > 2$ to $N = 2$ by ordering the parties in descending order by their
party votes and combining all parties except the first into one party $p'$.
Then one could assign the first party its constituencies according to the above procedure.
In the next step, one splits $p'$ into $p_2$ and $p''$ and repeats the procedure with the
constituencies not yet assigned. This is done repeatedly until all constituencies
are assigned. This procedure only has the disadvantage that the
party strongest in party votes can take away the relatively best constituencies from the other parties.
In the following example, the small party $p_3$ has its best result precisely where the large party $p_1$
also has its best result. But $p_1$, unlike $p_3$, has a \enquote{thick belly}. Even
if it were to give its strongest constituencies to $p_3$, it is still
strongly represented in the other constituencies, which is not the case with $p_3$. There the relative
vote shares go down pretty quickly with increasing $i$. Therefore, $p_3$ should get the first
constituencies and not $p_1$, even though it is the strongest force there.

\bigskip
\begin{center}
\includegraphics[width=0.8\textwidth]{diagrams/thick_belly.pdf}
\end{center}
\bigskip

We need a measure for the \enquote{relative importance} of a constituency for a party
that can be compared across parties. For this, we can determine from the matrix $R$ the matrix
$\Psi(R)$, for whose elements $\psi_{ij}$ should hold:

\[
  \psi_{ij} = (M-1)\,\frac{r_{ij}}{\displaystyle\sum_{\substack{i' \neq j,\, i'=1}}^{M} r_{i'j}}
  \quad \text{for} \quad \sum_{\substack{i \neq j,\, i=1}}^{M} r_{ij} \neq 0
  \quad \text{and $K$ otherwise,}
\]

where $K$ should be an arbitrarily chosen very large number ($K \gg \psi_{ij}\ \forall\,i,j$).

$\psi_{ij}$ gives the ratio of the relative vote share in constituency $i$ for party~$j$
to the sum of its relative vote shares in the other constituencies. Theoretically, this
sum could also be 0 if party~$j$ actually received no
votes in all constituencies except the $i$-th. Since each party must run in every constituency, this is very
unlikely. For this unlikely case, one can then simply set a very large number
$K$ that is assigned to the element instead of infinity. The more important a constituency is for
a party compared to all other constituencies, the larger $\psi_{ij}$ becomes. The
factor $(M-1)$ merely represents normalization so that the numerical values of $\Psi$ are more
human-readable. Through this normalization, $\Psi(R)$ is everywhere 1 with uniform distribution.

To now assign each constituency to a party, one first searches $\Psi$ for the largest element
$\psi_{ij}$. If there are several equally large $\psi_{ij}$, one randomly selects one of them.
The associated constituency $w_i$ is assigned to party $p_j$. The number of mandates $m_j$ still to be assigned to this party
is decreased by one and row $w_i$ is marked as no longer selectable.

Now one repeats the procedure with the still selectable $\psi_{ij}$ and the parties that
still have mandates to be assigned $m_j > 0$. Columns of parties with $m_j = 0$ are marked as not
selectable. The procedure is finished when $m_j = 0$ for all parties. I.e.,
no column is selectable anymore.

Thus, only the number of mandates determined by the party vote result per party is
distributed. The distribution starts with the regionally most unevenly distributed parties,
because they most urgently need certain constituencies.
More evenly distributed parties are not so dependent on a particular constituency.
They have performed equally well or equally poorly everywhere.

\runinheadul{Numerical Example:}

3 parties $p_1 \ldots p_3$ run in 5 constituencies $w_1 \ldots w_5$. For simplicity,
all constituencies should have the same size $G = 1000$. The following table lists the
party votes cast for the respective parties in the respective constituency:

\bigskip
\begin{center}
\begin{tabular}{l r r r}
  \toprule
           & $p_1$ & $p_2$ & $p_3$ \\
  \midrule
  $w_1$    & 650   & 100   & 250   \\
  $w_2$    & 630   & 200   & 170   \\
  $w_3$    & 600   & 280   & 120   \\
  $w_4$    & 550   & 350   & 100   \\
  $w_5$    & 450   & 500   & 50    \\
  \midrule
  $\Sigma$ & 2880  & 1430  & 690   \\
  \bottomrule
\end{tabular}
\end{center}
\bigskip

There are five constituencies, so 5 mandates must somehow be divided among
3 parties based on the election results. In principle, one must determine the relative vote shares $P_j$ of the
parties for the entire electoral area, multiply by the number of mandates,
and then round appropriately. Thus:
\[
  P_j = \frac{1}{S}\sum_{i=1}^{M} s_{ij}
  \quad \text{with} \quad S = \sum_{i=1}^{M}\sum_{j=1}^{N} s_{ij}
\]
\[
  m_j = \textit{round\_appropriately}(P_j \cdot M)
\]
But what does \enquote{round appropriately} mean? With simple standard rounding, the sum of the
$m_j$ does not necessarily equal $M$. More or fewer mandates can result from rounding than
are to be distributed.

Therefore, various procedures have been developed to get from a vote distribution to a
mandate distribution. The federal election uses the Sainte-Laguë/Schepers procedure.\footnote{\url{https://de.wikipedia.org/wiki/Sainte-Lagüe-Verfahren},
\url{https://www.bundeswahlleiterin.de/service/glossar/s/sainte-laque-schepers.html}}

The Sainte-Laguë/Schepers procedure yields for the above example:
$m_1 = 3$, $m_2 = 1$ and $m_3 = 1$.

Thus, 3 constituencies go to party $p_1$ and one constituency each to party $p_2$ and $p_3$. To
determine which constituencies go to which party, we first set up $\underline{R}$.
Since in our example the constituency size is constant at 1000, this is simply the content
of the above table divided by 1000:

\[
  \underline{R} = \begin{pmatrix}
    0.65 & 0.10 & 0.25 \\
    0.63 & 0.20 & 0.17 \\
    0.60 & 0.28 & 0.12 \\
    0.55 & 0.35 & 0.10 \\
    0.45 & 0.50 & 0.05
  \end{pmatrix}
\]

From this we determine $\Psi$ as

\[
  \underline{\Psi} = \begin{pmatrix}
    1.1659 & 0.3008 & 2.2727 \\
    1.1200 & 0.6504 & 1.3077 \\
    1.0526 & 0.9739 & 0.8421 \\
    0.9442 & 1.2963 & 0.6780 \\
    0.7407 & 2.1505 & 0.3125
  \end{pmatrix}
\]

The highest value of $\underline{\Psi}$ is $\psi_{13} = 2.2727$. I.e., constituency $w_1$
goes to party $p_3$, provided this party is still entitled to votes, i.e., provided $m_3 > 0$. That
is the case, so $w_1 \longrightarrow p_3$, $m_3$ is reduced by one, thus $m'_3 = 0$.

I.e., party $p_3$ gets no other constituency assigned and row 1 can no longer be selected from
$\underline{\Psi}$.

The next largest element in $\underline{\Psi}$ is $\psi_{52} = 2.1505$. Therefore, constituency
$w_5$ is assigned to party $p_2$ and $m_2$ is decreased by one. We have so far:
$w_1 \longrightarrow p_3$, $w_5 \longrightarrow p_2$, $m''_1 = 3$, $m''_2 = 0$, $m''_3 = 0$.

In the matrix, now only elements from the first column are selectable, all of which go to $p_1$,
which so far had not received any constituency at all. $p_1$ also still has enough
mandates, as we always distribute exactly as many mandates as constituencies with this procedure.

\runinheadul{Final Result:}

\begin{itemize}[label={}]
  \item $w_1 \longrightarrow p_3$
  \item $w_2 \longrightarrow p_1$
  \item $w_3 \longrightarrow p_1$
  \item $w_4 \longrightarrow p_1$
  \item $w_5 \longrightarrow p_2$
\end{itemize}

\subsection{Hungarian Algorithm for Optimizing the Assignment Problem}

The algorithm proposed above for assigning parties to constituencies has a problem: It
is \enquote{greedy}. It always takes the next largest element and thereby sometimes overlooks
more favorable combinations.

Constructed example with unnaturally high importance values, but shows the problem:

With $m_1 = 2$, $m_2 = 1$, $m_3 = 1$:

\[
  \underline{\Psi} = \begin{pmatrix}
    5 & 1 & 6 \\
    7 & 5 & 1 \\
    1 & 5 & 4 \\
    1 & 1 & 5
  \end{pmatrix}
\]

The greedy algorithm would proceed as follows:
\begin{itemize}[label={}]
  \item $w_2 \longrightarrow p_1$
  \item $w_1 \longrightarrow p_3$
  \item $w_3 \longrightarrow p_2$
  \item $w_4 \longrightarrow p_1$
\end{itemize}

The sum of importances would be: $\mathrm{IS}_{\mathrm{GREEDY}} = 7 + 6 + 5 + 1 = 19$

The optimal solution (in the sense of maximizing the assigned total importance) would however be:
\begin{itemize}[label={}]
  \item $w_1 \longrightarrow p_1$
  \item $w_2 \longrightarrow p_1$
  \item $w_3 \longrightarrow p_2$
  \item $w_4 \longrightarrow p_3$
\end{itemize}

$\mathrm{IS}_{\mathrm{OPTIMAL}} = 5 + 7 + 5 + 5 = 22$

The problem is that the greedy algorithm is not forward-looking. Through the assignment
$w_1 \to p_3$ (value 6 instead of 5 for $p_1$), it gains 1 point in the short term but loses
four points in row 4 (value 1 instead of 5 for $p_1$). Therefore, the sum of importances in the greedy
algorithm in this example is $4 - 1 = 3$ worse than the optimal solution.

There has long been a standard algorithm for such assignment problems: The
so-called \href{https://de.wikipedia.org/wiki/Ungarische\_Methode}{Hungarian Algorithm}
guarantees to find the optimal assignment in the sense of maximizing the sum of the assigned
relative importances. A version of this algorithm is also used in the simulation
by default.

\runinhead{Note:} Theoretically, it may happen that there is no global maximum, but
several equivalent assignments. In this case, the Hungarian algorithm returns any
one of them. In real elections at the federal or state level, this case is very
unlikely, as with many votes cast, the results per constituency and party will
always differ slightly.

The greedy algorithm with optimized implementation is not as computationally intensive as the
Hungarian algorithm.\footnote{Algorithmic order: Greedy algorithm: naive
implementation $O(M^2 N)$, optimized $O(M \cdot N \cdot \log(M \cdot N))$, Hungarian
algorithm: $O(M^2 N)$.} However, the difference plays practically no role in real elections.
Both algorithms deliver results on ordinary laptops in milliseconds.

Even though the greedy algorithm usually delivers the same results as the Hungarian algorithm,
I would recommend the Hungarian algorithm for such a sensitive topic as
constituency representation, since runtime plays practically no role. It is guaranteed that an
optimal assignment is found.
